%=====================================================================
%  Chapter 6 — Kinematics and Gait Planning
%=====================================================================
\chapter{Kinematics and Gait Planning}
\label{chap:kinematics}

\section{Introduction}
The software architecture of the previous chapter provides the channels
through which joint targets reach the robot. This chapter looks inside the
algorithm layer that \emph{computes} those targets. It presents, in turn,
the kinematic model of the legs, the inverse kinematics that turns desired
foot positions into joint angles, the inertial estimation and balance
control that keep the robot upright, and finally the Zero-Moment-Point
preview-control pipeline that generates a walking motion. The chapter ends
with an honest assessment of how far the resulting walk goes.

\section{Leg Kinematics}
\label{sec:fk}
Each leg is a chain of five revolute joints: two at the hip (roll and
pitch), one at the knee (pitch), and two at the ankle (pitch and roll).
\emph{Forward kinematics} --- composing these joint rotations along the
thigh and shank links --- gives the position of the foot relative to the
hip. For locomotion, however, the useful question is the inverse one:
given where the foot \emph{should} be, what joint angles place it there?
Because the leg has a simple, decoupled structure, this problem admits a
closed-form solution that separates cleanly into the sagittal plane
(forward/backward motion, Figure~\ref{fig:iksag}) and the frontal plane
(lateral motion, Figure~\ref{fig:ikfront}).

\begin{figure}[ht]
  \centering
  \begin{subfigure}{0.52\textwidth}
    \centering
    \includegraphics[height=6.2cm]{fig_ik_sagittal.png}
    \caption{Sagittal plane}
    \label{fig:iksag}
  \end{subfigure}
  \hfill
  \begin{subfigure}{0.38\textwidth}
    \centering
    \includegraphics[height=6.2cm]{fig_ik_frontal.png}
    \caption{Frontal plane}
    \label{fig:ikfront}
  \end{subfigure}
  \caption{Geometric inverse kinematics of the leg, decomposed into the
           sagittal plane (knee and hip/ankle pitch) and the frontal plane
           (hip and ankle roll).}
  \label{fig:ik}
\end{figure}

\section{Inverse Kinematics}
\label{sec:ik}
The robot's inverse kinematics is implemented as a \emph{closed-form},
two-link geometric solver. An earlier version used an iterative
damped-least-squares method, but it proved fragile: depending on its
starting guess it could converge to either a slightly-bent or a
near-straight leg, so the same foot target produced different poses
between runs. The closed-form solver removes this ambiguity --- given the
hip and foot positions it always returns the same human-like,
knee-forward solution --- which is essential for a repeatable gait.

The solver works as follows. In the frontal plane, the hip-roll angle is
the angle between the vertical and the line from hip to foot
($\arctan$ of the lateral over the vertical offset), and the ankle roll is
set to its opposite so that the sole stays flat on the ground. In the
sagittal plane, the straight-line distance from hip to foot together with
the known thigh and shank lengths fixes the knee angle by the
\emph{law of cosines}; the hip-pitch angle is the direction to the foot
minus the offset of the thigh from that line; and the ankle pitch is
chosen so that the hip, knee and ankle pitches sum to zero, again keeping
the foot flat. Two safeguards protect the solution: the hip-to-foot
distance is clamped just below full leg extension, and a minimum knee bend
is enforced, both of which keep the leg away from the singular,
fully-stretched configuration where the joint angles would be
ill-defined. A small forward offset is finally applied to the hip target,
because in a knee-bend the true body centre of mass sits slightly ahead of
the hip; shifting the target back keeps the centre of mass above the feet.

\section{Inertial Estimation}
\label{sec:imu-filter}
Balancing requires a clean estimate of how far the robot is tilted from
vertical, but neither inertial signal gives this directly: the
accelerometer measures the gravity direction but is corrupted by motion
noise, while the gyroscope measures rotation rate cleanly but drifts when
integrated over time. A \textbf{complementary filter} fuses the two by
trusting the integrated gyroscope over the short term and gently
correcting it toward the accelerometer estimate over the long term. With
its coefficient set to $0.98$, the filter follows fast motions smoothly
yet remains free of long-term drift. It consumes the raw inertial topic
(from the simulated sensor or, on the real robot, from the ESP32) and
publishes the filtered pitch and roll used by the balance controller.

\section{Balance Control}
\label{sec:balance}
Standing balance is maintained by an inner feedback loop that adjusts the
ankle joints in reaction to the measured tilt. Rather than three
independent PID gains per axis, the controller uses a \textbf{Linear-
Quadratic Regulator (LQR)}. The motivation is physical: the ankle couples
the pitch and roll axes, so tuning them independently is suboptimal. The
LQR instead treats the full state --- pitch, roll, and their rates --- and
computes a single optimal gain matrix offline by solving the algebraic
Riccati equation, trading off a penalty on tilt error against a penalty on
ankle effort. At runtime the correction is a single matrix multiplication.
Following the same linear-inverted-pendulum model used for gait
(Section~\ref{sec:gait}), the controller takes the nominal joint angles
produced by the kinematics, adds the computed ankle correction, and
republishes the corrected command, leaving the rest of the system
unchanged.

\section{Gait Generation by ZMP Preview Control}
\label{sec:gait}
Walking is generated using the \textbf{Zero-Moment-Point preview-control}
method of Kajita (2003)~\cite{kajita2003}, built on the Linear Inverted
Pendulum Model (LIPM). The LIPM abstracts the whole robot as a point mass at constant
height: under this model the ZMP --- the point introduced in
Chapter~\ref{chap:walking} --- depends on the centre-of-mass position and
its acceleration. The walking problem then becomes: find a centre-of-mass
trajectory whose resulting ZMP stays inside the support foot at all times.
Figure~\ref{fig:gaitcycle} shows the relationship between the foot
placement, the centre-of-mass path and the ZMP across the
single- and double-support phases.

\begin{figure}[ht]
  \centering
  \includegraphics[width=0.9\textwidth]{fig_gait_cycle.png}
  \caption{Walking in the sagittal (a) and frontal (b) planes: the
           centre-of-mass trajectory and the ZMP evolve across the
           double-support (DS) and single-support (SS) phases.}
  \label{fig:gaitcycle}
\end{figure}

The key difficulty is that the ZMP at the present instant depends on the
\emph{future} motion of the centre of mass. The preview controller solves
this by looking ahead~\cite{katayama1985}: it is an optimal regulator,
derived from the same Riccati machinery as the balance controller,
augmented with a
\emph{preview} term that anticipates the reference ZMP over a horizon of
about one and a half seconds. This lets the centre of mass start shifting
toward the next footstep \emph{before} the foot is placed, exactly as a
human leans into a step. The method has since been extended with auxiliary
ZMP control and more general preview formulations~\cite{kajita2006,park2007}.

The generation is organised as a three-stage pipeline:
\begin{enumerate}[itemsep=2pt,topsep=4pt]
  \item a \textbf{footstep planner} lays out a timed sequence of left/right
        foot placements along the desired path, with shorter first steps to
        ramp up from standing;
  \item a \textbf{ZMP reference} stage converts those footsteps into a
        target ZMP over time --- under the stance foot during single
        support, interpolating between feet during double support;
  \item the \textbf{preview controller} turns this ZMP reference into a
        smooth centre-of-mass trajectory.
\end{enumerate}
In parallel, the swing foot follows a lifted trajectory between its
successive placements (Figure~\ref{fig:wpg}). The centre-of-mass and foot
trajectories are then handed to the inverse kinematics of
Section~\ref{sec:ik}, which produces the joint angles, and the balance
controller adds its ankle correction on top. To react to disturbances, the
pipeline also closes a \emph{capture-point} loop: from the measured tilt
and tilt rate it estimates where the robot would have to step to come to
rest, and steers the swing foot's landing toward that point.

\begin{figure}[ht]
  \centering
  \includegraphics[width=0.7\textwidth]{fig_gait_wpg.png}
  \caption{Swing-foot trajectory over the planned footstep sequence
           produced by the walking-pattern generator.}
  \label{fig:wpg}
\end{figure}

\section{The Complete Motion Pipeline}
\label{sec:pipeline}
Figure~\ref{fig:motionpipe} assembles these components into the full
motion pipeline. The feed-forward gait chain (footstep planner $\to$ ZMP
reference $\to$ preview controller) produces the centre-of-mass and foot
trajectories; the inverse kinematics converts them to nominal joint
angles; and the balance controller adds an ankle correction driven by the
inertial feedback. The same inertial feedback feeds the capture-point
adjustment, closing the loop back onto the foot placement.

\begin{figure}[ht]
  \centering
  \begin{tikzpicture}[
      font=\footnotesize, >=Latex, node distance=6mm and 9mm,
      ff/.style={draw, rounded corners=2pt, align=center, thick,
                 fill=blue!7, minimum height=8mm, text width=20mm},
      ik/.style={draw, rounded corners=2pt, align=center, thick,
                 fill=orange!12, minimum height=8mm, text width=20mm},
      fb/.style={draw, rounded corners=2pt, align=center, thick,
                 fill=green!10, minimum height=8mm, text width=20mm},
      dst/.style={draw, rounded corners=2pt, align=center, thick,
                  fill=gray!12, minimum height=8mm, text width=18mm},
    ]
    \node[ff] (fp) {Footstep planner};
    \node[ff, right=of fp] (zr) {ZMP reference};
    \node[ff, right=of zr] (pc) {Preview controller (LIPM)};
    \node[ik, right=of pc] (ik) {Inverse kinematics};
    \node[ik, right=of ik] (bal) {Balance (LQR)};
    \node[dst, right=of bal] (rob) {Robot / sim};

    \node[fb, below=10mm of bal] (imu) {IMU $\to$ filter $\to$ tilt};

    \draw[->] (fp) -- (zr);
    \draw[->] (zr) -- (pc);
    \draw[->] (pc) -- node[above,font=\tiny]{CoM} (ik);
    \draw[->] (ik) -- node[above,font=\tiny]{$\theta$} (bal);
    \draw[->] (bal) -- (rob);
    % foot trajectory branch
    \draw[->] (fp.south) |- ($(pc.south)+(0,-6mm)$) -| node[pos=0.25,below,font=\tiny]{foot traj.} (ik.south);
    % feedback
    \draw[->] (rob.south) |- (imu.east);
    \draw[->] (imu) -- node[right,font=\tiny]{$\Delta\theta$ ankle} (bal);
    \draw[->] (imu.west) -| node[pos=0.25,below,font=\tiny]{capture point} (fp.south);
  \end{tikzpicture}
  \caption{The complete motion pipeline: a feed-forward ZMP gait chain
           feeds the inverse kinematics; the balance controller and the
           capture-point adjustment close the loop with inertial feedback.}
  \label{fig:motionpipe}
\end{figure}

\section{Results and Limitations}
\label{sec:gait-results}
The kinematics and standing balance work well: the closed-form inverse
kinematics is deterministic and produces correct, repeatable leg poses,
and with the balance controller active the robot holds a deep crouch
stably for extended periods in simulation. The gait pipeline produces
correct ZMP references, centre-of-mass trajectories and alternating foot
motions, and on good runs the robot \emph{marches in place} for several
balanced steps.

Robust forward walking, however, remains the open frontier of the project.
The fundamental limitation is that the robot is driven by stiff
position-controlled hobby servos with no torque sensing, and the lateral
balance during single support is only marginally stable: a small
side-to-side oscillation tends to grow over several steps. This is a
known, method-level limitation rather than a tuning issue --- closing it
robustly calls for full capture-point step-adjustment feedback or a
learned controller (reinforcement learning), which is identified as
post-project work in Chapter~\ref{chap:conclusion}. The detailed
experimental results are presented in Chapter~\ref{chap:validation}.

\section{Conclusion}
This chapter presented the algorithms that compute the robot's joint
targets: a closed-form leg inverse kinematics, a complementary filter for
inertial estimation, an LQR balance controller, and a Kajita-style ZMP
preview-control gait pipeline closed with capture-point feedback. Together
they let the robot stand, balance and step in simulation, with robust
forward walking remaining the principal open challenge. The next chapter
turns to the highest layer of the system: the artificial-intelligence
command interface that lets a user drive these motions in natural
language.
